\section{Edge AI and TinyML Systems}
\label{sec:edge-ai-stm32}

Edge Artificial Intelligence (Edge AI) refers to the deployment of machine learning models directly on embedded devices, such as microcontrollers (MCUs) and digital signal processors (DSPs), rather than on centralized cloud servers. This shift originates from the demand for low latency, reduced bandwidth consumption, data privacy, and improved energy efficiency \cite{warden2019tinyml}. However, deploying sophisticated AI on resource-constrained hardware like the STM32 family presents precise technical challenges.

\subsection{Challenges of TinyML}
The field of Tiny Machine Learning (TinyML) focuses on enabling inference on devices with power consumption typically in the milliwatt range \cite{warden2019tinyml}. The primary constraints include:

\begin{itemize}
    \item \textbf{Memory Limitations}: MCUs generally have restricted RAM (hundreds of KB) and Flash memory (a few MB). This structural limit requires hardware-aware model compression techniques like quantization and pruning \cite{st_cubeai_docs}.
    \item \textbf{Computational Power}: Without high-end GPUs, models must be explicitly optimized for execution on ARM Cortex-M cores or dedicated hardware accelerators, such as the STM32 NPU.
    \item \textbf{Energy Efficiency}: Many embedded applications are battery-powered, demanding architectures that minimize active clock cycles and maximize sleep states.
\end{itemize}

\subsection{TinyMLOps: MLOps for Resource-Constrained Devices}

The rapid expansion of Artificial Intelligence on embedded systems forced the adaptation of traditional Machine Learning Operations (MLOps) principles to the specific constraints of micro-controllers. This domain, known as \textbf{TinyMLOps}, addresses the lifecycle management of AI models in environments where memory, compute, and power are severely limited \cite{warden2019tinyml}. TinyMLOps provides the methods and the tooling necessary to face the challenges outlined above.

\subsubsection{Distinction from Cloud-Scale MLOps}

While traditional MLOps manages model deployment across cloud clusters, massive datasets, and microservices, TinyMLOps prioritizes strict hardware awareness. The main differences are summarized in Table~\ref{tab:mlops_vs_tinymlops}.

\begin{table}[h]
    \centering
    \caption{Comparison between traditional Cloud MLOps and TinyMLOps \cite{mlops_productivity_2024}.}
    \label{tab:mlops_vs_tinymlops}
    \begin{tabular}{|l|l|l|}
        \hline
        \textbf{Feature} & \textbf{Cloud MLOps}        & \textbf{TinyMLOps}              \\
        \hline
        Target Hardware  & GPU Clusters / Data Centers & MCUs (KB-scale RAM)             \\
        Optimization     & Throughput / Latency        & Memory / Energy Efficiency      \\
        Deployment Unit  & Docker Containers / APIs    & Firmware Binary Images          \\
        Connectivity     & Persistent / High-bandwidth & Intermittent or Offline         \\
        Metric Priority  & Prediction Accuracy         & Accuracy vs. Resource Footprint \\
        \hline
    \end{tabular}
\end{table}

\subsubsection{Core Pillars of TinyMLOps}

A systematic TinyMLOps framework typically implements three fundamental pillars:

\begin{itemize}
    \item \textbf{Hardware-Aware Optimization}: Unlike cloud models that scale easily with more VRAM, embedded models undergo structural transformations (e.g., quantization, pruning, and "neural surgery") to fit into specific MCU memory profiles.

    \item \textbf{Integrated Automation ("Plumbing")}: TinyMLOps requires the direct integration of traditionally disjoint toolchains. It connects Python-based training environments, C-based firmware generators (like STM32Cube), and binary flashing utilities.

    \item \textbf{Reproducibility}: Defining the end-to-end workflow as a declarative graph (e.g., via LangGraph) ensures that engineers can reconstruct and verify an optimized model consistently across different hardware revisions \cite{wang2024agentailanggraphmodular}.
\end{itemize}

Adopting a TinyMLOps approach transitions the development of embedded AI from an expert-intensive manual process to an automated engineering discipline.

\subsection{STM32 Microcontrollers Overview}

To fully appreciate the scope of TinyMLOps, it is important to define the target hardware ecosystem. The STM32 portfolio, built upon ARM Cortex-M processors (such as Cortex-M4, Cortex-M33, and Cortex-M7), is divided into specific families designed for different application domains. More information on all families is given in appendix A.

The main categories relevant to Edge AI deployment include:

\begin{itemize}
    \item \textbf{Foundation Family (e.g., STM32F4)}: These general-purpose microcontrollers feature RAM between 96 and 128 KB and Flash memory of approximately 500 to 512 KB. They power applications involving basic audio processing or simple sensor data analysis.
    \item \textbf{Low Power Family (e.g., STM32L4)}: Destined for applications where energy consumption is critical, this family offers extended resources, with RAM scaling up to 640 KB and Flash reaching 2 MB. They are highly suitable for low-power vision or continuous monitoring tasks.
    \item \textbf{High Performance Family (e.g., STM32H7)}: When the application relies strictly on the core's high performance, the H7 family delivers. With RAM capacities up to 1.4 MB, it hosts heavier models, including complex object detection computer vision architectures.
\end{itemize}

\subsection{Model-Hardware Compatibility Matrix}
To provide concrete guidance on model selection for these resource-constrained devices, Table~\ref{tab:model_board_compat_summary} presents a condensed compatibility analysis of vision models from the STM32 Model Zoo across common MCU families. These models target \textbf{image classification} and \textbf{pattern recognition} tasks, ranging from simple digit recognition (ST MNIST) to object categorization (MobileNet and ResNet variants). The analysis considers both uncompressed (FP32) and INT8-quantized versions, demonstrating the effect of compression in enabling computer vision applications on lower-tier devices.

\begin{table}[h!]
    \centering
    \caption{Representative Model-Board Compatibility (INT8 Quantized)}
    \label{tab:model_board_compat_summary}
    \resizebox{\textwidth}{!}{%
        \begin{tabular}{|l|c|c|c|c|c|c|c|}
            \hline
            \textbf{Model}       & \textbf{Input} & \textbf{RAM}    & \textbf{Flash}  & \textbf{F401}       & \textbf{L476}      & \textbf{H743}    & \textbf{N6570}   \\
                                 & \textbf{Size}  & \textbf{(INT8)} & \textbf{(INT8)} & \textbf{64KB/256KB} & \textbf{128KB/1MB} & \textbf{1MB/2MB} & \textbf{5MB/4MB} \\
            \hline
            ST MNIST             & 28×28          & 5 KB            & 25 KB           & \checkmark          & \checkmark         & \checkmark       & \checkmark       \\
            FD-MobileNet 0.25    & 96×96          & 80 KB           & 180 KB          & $\times$            & \checkmark         & \checkmark       & \checkmark       \\
            MobileNet V2 0.35    & 128×128        & 150 KB          & 280 KB          & $\times$            & \textasciitilde    & \checkmark       & \checkmark       \\
            ResNet V1 32 (CIFAR) & 32×32          & 200 KB          & 450 KB          & $\times$            & $\times$           & \checkmark       & \checkmark       \\
            MobileNet V2 1.0     & 224×224        & 600 KB          & 1.2 MB          & $\times$            & $\times$           & \textasciitilde  & \checkmark       \\
            \hline
        \end{tabular}%
    }
    \vspace{2mm}
    \begin{flushleft}
        \small
        Legend: \checkmark = Fits comfortably ($<$80\% resources); \textasciitilde = Tight fit (80-100\%); $\times$ = Does not fit
    \end{flushleft}
\end{table}

We derived key observations from this analysis:
\begin{itemize}
    \item \textbf{INT8 Quantization Impact}: Compression reduces the model footprint by approximately 4$\times$ compared to FP32. It enables the deployment of models like MobileNet V2 0.35 on mid-tier MCUs (L476) that would otherwise exceed hardware limits \cite{st_cubeai_docs}.
    \item \textbf{Vision Models on F401}: Standard image classification architectures remain incompatible with the STM32F401 (64KB RAM) even with extreme quantization, limiting this platform to audio-based inference or ultra-lightweight custom Neural Networks.
    \item \textbf{H7 Family as Sweet Spot}: The STM32H743/H747 (1MB RAM) emerges as the optimal target for production-grade vision AI, executing MobileNet variants at 224$\times$224 resolution with adequate latency.
\end{itemize}

The complete compatibility matrix, including FP32 resource requirements and additional model architectures (EfficientNet, SqueezeNet, ResNet50), is available in Appendix~\ref{app:model_compat_full}.

\subsection{The STM32 Developer Toolchain}
STMicroelectronics maintains an ecosystem for migrating high-level AI models to optimized C code. The core tools integrated within our workflow include:

\begin{itemize}
    \item \textbf{STM32CubeMX}: A graphical tool for configuring MCU peripherals and generating initialization code. In an automated MLOps workflow, this tool provides the "hardware context" (e.g., pinout, clock settings) necessary for firmware integration.
    \item \textbf{STEdgeAI / X-CUBE-AI}: This specialized utility converts pre-trained models (from TensorFlow Lite, Keras, or ONNX) into optimized C code for STM32. It yields critical feedback on memory footprint (RAM/Flash) and execution cycles, acting as the interface between the AI and the embedded domains \cite{st_cubeai_docs}.
\end{itemize}

The orchestration of these tools typically forces engineers into iterative manual cycles. Developers repeatedly adjust model parameters in Python, execute the conversion tool, and supervise hardware compatibility. The "GUI-centric" operation of tools like STM32CubeMX acts as a bottleneck for industrial MLOps pipelines because it prevents version control and automation within continuous integration systems. To resolve this, the framework introduced in this thesis automates the "plumbing" of the pipeline via headless agentic workflows.
